% ===========================================================================
%  第 11 章 结论
% ===========================================================================
\section{结论}

\subsection{工作总结}

本报告对纵列双发矢量推力飞行器飞控固件中的 EKF 组合导航系统完成了从理论到
实现的完整分析：沿``坐标系与随机过程基础 $\rightarrow$ 地球模型 $\rightarrow$
捷联编排 $\rightarrow$ 误差状态线性化 $\rightarrow$ 量测模型 $\rightarrow$
鲁棒性机制 $\rightarrow$ 任务集成 $\rightarrow$ 仿真验证''的技术主线，将
3968 行滤波器核心与 1769 行导航任务中的每一处矩阵块、增益调度与门限常量还原
为严格的数学推导，并基于固件真实参数构建独立数值仿真对关键行为进行了定量
验证。

\subsection{主要结论}

\begin{enumerate}
  \item \textbf{滤波器架构}：固件采用 15 状态误差状态 EKF
    $\left[\delta\vp{p}^n,\ \delta\vp{v}^n,\ \delta\vp{\beta}^b,\
    \delta\vp{b}_a,\ \delta\vp{b}_g\right]$，姿态误差为体系右乘
    （$q = q_{nom}\otimes\Exp(\delta\vp{\beta}^b)$）。该定义使
    $\m{F}_{\beta v} = -\m{C}_n^b\m{M}$（第 5.3.9 节），与导航系左乘定义
    差一个旋转，数学上同构但线性化点选择影响大修正时的精度——固件以指数映射
    注入（第 5.7 节）处理该问题。

  \item \textbf{机械编排}：双子样圆锥/划摇补偿（$\frac23\Delta\vp{\theta}_1
    \cross\Delta\vp{\theta}_2$ 与流式单子样版本）提供 $dt^3$ 阶增量精度；
    中点重力/科氏力与速度中点位置积分使姿态、速度、位置离散化均达二阶；
    导航系转动补偿与静止初始化时的地球自转零偏分离（第 4.8 节）保证长时间
    静置无慢漂移。

  \item \textbf{过程噪声}：Simpson 离散化
    $Q_d = \frac{dt}{6}\left(Q_c + 4\Phi_{\frac12}Q_c\Phi_{\frac12}\T +
    \Phi Q_c\Phi\T\right)$（第 2.6 节）比零阶保持保留二阶耦合；正交不变性
    将结构化 $Q_c$ 化简为对角直写，数学等价而计算廉价。

  \item \textbf{鲁棒性纵深}：NIS 门控（统计层）、物理残差门限（语义层）、
    失联协方差膨胀 + 低增益重捕获（恢复层）、GNSS 延迟回放重传播（时间对齐
    层）与输入合法性防御（数据层）构成五层防御。第 10 章仿真分别验证了各层
    在对应场景下的有效性。

  \item \textbf{可观性边界}：纯静止下水平加速度零偏不可观（与倾角组合可观），
    垂直零偏与三轴陀螺零偏可观；无磁/无 GNSS 时 yaw 不可观。这些边界直接
    决定了固件参数设计（如初始航向协方差 $\pi$、陀螺零偏 Markov 驱动
    0.3\,dps）与使用方式（无 GNSS 时仅姿态可用）。

  \item \textbf{任务集成}：200\,Hz 主任务通过``统一静止检测 + 三档辅助调度 +
    条件性 AHRS 辅助 + GNSS 动态权重''在可用性与精度之间自适应；独立垂直/
    水平 KF 提供冗余备份而不污染主链路；EKF/AHRS/DETA100 三路输出切换通过
    航向偏置慢校正保证连续性。
\end{enumerate}

\subsection{仿真验证的意义}

第 10 章的独立数值仿真（约 800 行 Python，固件真实参数）完成了对滤波器行为
的端到端验证：稳态融合精度与协方差一致性（S2）、失联漂移与膨胀覆盖（S3）、
延迟回放的时间对齐收益（S4）、双矢量航向的漂移消除（S5）、NIS 门控的野值
拦截（S6）以及静止辅助的可观性边界（S1）。仿真结果与理论推导一致，为固件
参数调优与使用边界提供了定量依据。

\subsection{局限与后续工作}

\begin{enumerate}
  \item \textbf{仿真保真度}：当前仿真以高斯白噪声建模传感器误差，未包含真实
    共轴双桨震动谱、时变零偏、温度漂移与安装误差。后续可引入实测 IMU 数据
    回放（黑盒数据驱动）以标定噪声模型并验证仿真器；
  \item \textbf{水平加速度零偏的可观化}：S1 表明水平零偏需 GNSS 运动激励或
    已知机动解耦。固件可考虑在 GNSS 有效时增加零偏慢速估计的显式调度，或
    引入 RTK 高精度速度观测加速解耦；
  \item \textbf{气压计状态化}：第 6.7 节指出当前气压高度量测无零偏状态，QNH
    偏差会与 GNSS 绝对高度冲突。后续可将 QNH/气压零偏增广为状态量，消除
    该限制；
  \item \textbf{延迟标定}：\code{GNSS\_DELAY\_S}=15\,ms 为无 PPS 的初始估计。
    接入 PPS/TIMEPULSE 或按 iTOW 与 MCU 时间实测标定后，可在
    platformio.ini 覆盖该宏（固件注释已预留），进一步收紧延迟回放精度；
  \item \textbf{与实机数据的闭环验证}：建议将实机飞行的 IMU/GNSS 数据回放
    至仿真器，对比 EKF 输出与实机遥测，验证仿真器与固件的逐帧一致性。
\end{enumerate}

\subsection{结语}

EKF 组合导航是飞控系统中``理论密度''最高的子系统之一：它把捷联惯导的几何
编排、随机过程的统计建模、卡尔曼滤波的估计理论、地球模型的物理常数与嵌入式
实现的数值防御熔于一炉。本报告以 3968 行固件代码为底本完成的这一份理论-
算法-实现-仿真四层映射，既是对现有系统的完整技术档案，也是后续演进
（RTK 融合、视觉辅助、零偏状态化）的数学起点。报告中所有公式均可与固件源码
逐行对照，所有仿真数据均可由随附脚本复现，为该系统的评审、传承与迭代提供了
坚实基础。
